% =====================================================================
%  corps/02-exemples-multilingues.tex
%  Chapitre vitrine + MODE D'EMPLOI : chaque fonctionnalité est présentée
%  avec son code source (environnement « latexcode ») puis son rendu.
% =====================================================================
\chapter{Exemples multilingues et mode d'emploi}
\label{chap:exemples}

Ce chapitre illustre l'insertion de texte dans différentes écritures et
sert en même temps de \emph{mode d'emploi} : pour chaque fonctionnalité,
le code source est affiché, suivi de son rendu. Chaque langue secondaire
est pré-configurée dans \code{config/langues.tex} avec sa police adéquate
(\code{config/polices.tex}).

% ---------------------------------------------------------------------
\section{Fragments courts en ligne}

Le modèle définit des commandes de confort pour insérer un mot ou une
expression dans une écriture donnée. On les utilise ainsi (remplacez le
texte entre accolades par votre contenu dans l'écriture voulue) :

\begin{latexcode}
Arabe    : \ar{ ... texte arabe ... }      % droite à gauche
Russe    : \ru{ ... texte russe ... }
Chinois  : \zh{ ... texte chinois ... }
Japonais : \ja{ ... texte japonais ... }
Grec     : \gr{ ... texte grec ... }
Tigrinya : \ti{ ... texte tigrinya ... }
\end{latexcode}

\noindent \textbf{Rendu :}

\begin{itemize}
  \item Arabe : \ar{مرحبا بالعالم}
  \item Russe : \ru{Здравствуй, мир}
  \item Chinois : \zh{你好，世界}
  \item Japonais : \ja{こんにちは世界}
  \item Grec : \gr{Γειά σου, κόσμε}
  \item Tigrinya : \ti{ሰላም ዓለም}
\end{itemize}

Ces fragments cohabitent sans difficulté dans un paragraphe français, y
compris pour une langue de droite à gauche. Par exemple, le code

\begin{latexcode}
comme l'arabe \ar{ ... } qui reprend ensuite le fil du texte.
\end{latexcode}

\noindent produit : comme l'arabe \ar{النص العربي} qui reprend ensuite le
fil du texte français.

% ---------------------------------------------------------------------
\section{Passages longs (environnements)}

Pour un passage plus long, on utilise l'environnement de la langue, qui
gère la césure et la direction du texte. La syntaxe générale est :

\begin{latexcode}
\begin{Arabic}
  ... votre paragraphe en arabe ...
\end{Arabic}

\begin{russian}
  ... votre paragraphe en russe ...
\end{russian}
\end{latexcode}

\begin{quote}\small
\textbf{Attention.} L'environnement arabe s'écrit avec une \emph{majuscule}
(\code{Arabic}) pour éviter un conflit avec la commande \code{\arabic} de
LaTeX. Les autres s'écrivent en minuscules (\code{russian}, \code{chinese},
\code{japanese}, \code{greek}).
\end{quote}

\subsection{Arabe (droite à gauche)}

Un paragraphe complet en arabe est justifié et aligné à droite, avec un
ordre des lignes et une césure corrects (paquet \code{bidi}) :

\begin{Arabic}
تُسَاعِدُ اللُّغَاتُ المُخْتَلِفَةُ عَلَى التَّوَاصُلِ بَيْنَ النَّاسِ وَالثَّقَافَاتِ. وَمَعَ تَطَوُّرِ التِّقْنِيَاتِ الرَّقْمِيَّةِ، أَصْبَحَ مِنَ السَّهْلِ كِتَابَةُ النُّصُوصِ وَقِرَاءَتُهَا بِأَحْرُفٍ وَاتِّجَاهَاتٍ مُخْتَلِفَةٍ. وَيُسَاعِدُ دَعْمُ النُّصُوصِ مِنَ الْيَمِينِ إِلَى الْيَسَارِ عَلَى عَرْضِ اللُّغَةِ الْعَرَبِيَّةِ بِشَكْلٍ وَاضِحٍ وَطَبِيعِيٍّ.
\end{Arabic}

\subsection{Russe (cyrillique)}
\begin{russian}
Русский язык использует кириллический алфавит. Этот абзац демонстрирует
корректный набор кириллического текста, включая переносы и
типографику.
\end{russian}

\subsection{Chinois}
\begin{chinese}
这是一段中文示例文字，用于演示本模板对中日韩表意文字的排版支持。
中文通常不使用空格分词。
\end{chinese}

\subsection{Japonais}
\begin{japanese}
これは日本語のサンプル段落です。ひらがな、カタカナ、漢字が
混在しても正しく組版されることを示します。
\end{japanese}

\subsection{Grec}
\begin{greek}
Αὕτη ἡ παράγραφος εἶναι γραμμένη στὰ ἑλληνικά, γιὰ νὰ δείξει τὴν
ὑποστήριξη τοῦ ἑλληνικοῦ ἀλφαβήτου, μαζὶ μὲ τόνους καὶ πνεύματα.
\end{greek}

% ---------------------------------------------------------------------
\section{Tableaux (booktabs)}

Les tableaux se composent avec \code{booktabs} (filets \code{\toprule},
\code{\midrule}, \code{\bottomrule}) et \code{tabularx} (colonne \code{X}
de largeur automatique). Ils acceptent les commandes multilingues :

\begin{latexcode}
\begin{table}[htbp]
  \centering
  \begin{tabularx}{\linewidth}{@{}l X@{}}
    \toprule
    \textbf{Langue} & \textbf{« traduction »} \\
    \midrule
    Français & traduction \\
    Arabe    & \ar{ ... } \\
    Chinois  & \zh{ ... } \\
    \bottomrule
  \end{tabularx}
  \caption{Légende du tableau.}
  \label{tab:traduction}
\end{table}
\end{latexcode}

\noindent \textbf{Rendu :}

\begin{table}[htbp]
  \centering
  \begin{tabularx}{\linewidth}{@{}l X@{}}
    \toprule
    \textbf{Langue} & \textbf{« Traduction » (le mot pour « traduction »)} \\
    \midrule
    Français  & traduction \\
    Arabe     & \ar{ترجمة} \\
    Russe     & \ru{перевод} \\
    Chinois   & \zh{翻译} \\
    Japonais  & \ja{翻訳} \\
    Grec      & \gr{μετάφραση} \\
    Tigrinya  & \ti{ትርጉም}\\
    \bottomrule
  \end{tabularx}
  \caption{Le terme « traduction » dans plusieurs langues et écritures.}
  \label{tab:traduction}
\end{table}

% ---------------------------------------------------------------------
\section{Notes de bas de page et citations}

Une note de bas de page s'insère avec \code{\footnote{...}} et peut, elle
aussi, contenir plusieurs écritures :

\begin{latexcode}
\dots{} plusieurs écritures.\footnote{Arabe \ar{ ... }, russe \ru{ ... }.}
\end{latexcode}

\noindent \textbf{Rendu :} une note peut contenir plusieurs
écritures.\footnote{Exemple : arabe \ar{هامش}, russe \ru{сноска},
chinois \zh{脚注}, japonais \ja{脚注}.}

Les références bibliographiques se citent avec \code{\autocite{clé}}
(entre parenthèses) ou \code{\textcite{clé}} (dans la phrase) :

\begin{latexcode}
D'après \textcite{barhudarov1975}, la traduction \dots{}  \textcite{xietianzhen1999, aljahiz}
\dots{} la traductologie outillée \autocite{schlechtweg-etal-2020-semeval}.
\end{latexcode}

\noindent \textbf{Rendu :} D'après \textcite{barhudarov1975}, la traduction \dots{} \textcite{xietianzhen1999, aljahiz} 
\dots{} la traductologie outillée \autocite{schlechtweg-etal-2020-semeval}.

Comme démontré ci-dessus, les références à des auteurs dont le nom s'écrit dans des caractères non latins s'affichent et se
trient correctement grâce à \code{biber}.

\code{biber} permet également d'ajouter des \enquote{pre-notes} et \enquote{post-notes} :
\begin{latexcode}
\dots{} nous pouvons mentionner \textcite[originalement du 9ème siècle,][qui traduit Aristote]{aljahiz}
\end{latexcode}

\noindent \textbf{Rendu :} \dots{} nous pouvons mentionner \textcite[originalement du 9ème siècle,][qui traduit Aristote]{aljahiz}